\begin{abstract}
Text-to-video generators still render garbled text on signs, screens and posters. We study \emph{repairing} such text: given a generated clip, a target string and a hint box, make the text read correctly while leaving everything else as it was. Generative video text editors produce natural-looking glyphs, but they frequently keep or invent wrong characters, rewrite text that was already correct, and re-render the background. We present \method, a training-free, correctness-first repair pipeline built around three evidence gates. A \emph{preservation gate} leaves a clip untouched when multi-frame OCR already reads the target; a \emph{localization gate} lays the new text out on the old glyph instance when the evidence is strong and falls back to the tracked box otherwise; and an \emph{appearance gate} enables a non-parametric source appearance transport (SAT) that carries glow and material from the old glyphs to the new ones only where it measurably helps. The target string itself is rendered deterministically, so its spelling cannot drift, and only pixels inside the text region change. On a benchmark of 526 generated clips spanning seven carrier materials in English and Chinese, \method reaches \rOursRead\% readability against \rSteerRead\% for the state-of-the-art generative editor \steer (\rOursNet{} vs.\ \rSteerNet{} net repairs; a pre-registered second reader from another model family confirms the advantage, \jIvlOurs{} vs.\ \jIvlSteer\%) and keeps the background intact (\rOursBG{} vs.\ \rSteerBG\,dB PSNR outside the text region). In a blind study, raters preferred SAT over plain pasting for material fidelity in 87\% of non-tied judgements. We also report, with pre-registered criteria where applicable, what did not work---a learned 1.3B refiner, glyph-only erasure, a carrier-fit check, reuse of correctly generated source glyphs, and baseline-aligned layout with glyph-matched fonts---and we characterise the appearance gap to generative editing that remains.
\end{abstract}
